\documentclass[10pt,a4paper]{amsart}
\usepackage[margin=19mm]{geometry}
\usepackage{amsmath,amssymb,mathtools}
\usepackage[scr=rsfs,cal=euler]{mathalfa}
\usepackage{xfrac}
\usepackage[unicode,hidelinks,bookmarks=false]{hyperref}
\newcommand{\ie}{i.e.\ }
\newcommand{\cf}{cf.\ }
\newcommand{\resp}{respectively\ }
\newcommand{\Subsection}{\subsection}
\providecommand{\nobreakdash}{\nobreak\hbox{-}}
\setlength{\emergencystretch}{2em}
\multlinegap0pt
\usepackage{bm}
\usepackage{bbm}
\usepackage{dsfont}
\usepackage{xspace}
\usepackage{cancel}
\usepackage{mathtools}
\usepackage{amsthm}
\makeatletter
\@ifundefined{lemma}{\newtheorem{lemma}{Lemma}}{}
\@ifundefined{theorem}{\newtheorem{theorem}{Theorem}}{}
\makeatother
\title[Efficient approximations of matrix multiplication using trun]{Efficient approximations of matrix multiplication using truncated decompositions}
\author{Suvendu Kar, Hariprasad M., Sai Gowri J. N.,, Murugesan Venkatapathi}
\date{Source: arXiv:2504.19308v4 (CC BY 4.0)}
\begin{document}
\maketitle

\noindent\emph{Source and licence.} Efficient approximations of matrix multiplication using truncated decompositions. Version arXiv:2504.19308v4 (CC BY 4.0), distributed under the Creative Commons Attribution 4.0 International licence, as recorded in the arXiv licence metadata for this exact version and shown on the abstract page https://arxiv.org/abs/2504.19308v4. Full rights notice: licenses/arxiv-2504.19308-CC-BY-4.0.txt.\par\smallskip

\noindent\textit{Editorial scope.} Counted unit: the Theorem whose source label is \texttt{th:cd-deterministic} in the article (source0.tex lines 757--766, source label \texttt{th:cd-deterministic}), delivered together with the article's own complete original proof of it (source0.tex lines 768--837, one \texttt{proof} environment of 70 lines). No mathematics is written, repaired or re-derived: statement and proof are the article's own text line for line. Required context retained uncounted: lem:circulant-conjugation (lines 709--712). Every definition, display and label of those lines is retained unchanged. Omitted from this package and not redistributed: 1--708, 713--756, 767--767, 838--1753. Nothing inside a retained excerpt is omitted or altered: every retained body is delivered exactly as the article's lines are written, so no omission receipt of any kind accompanies this package. Citations: the retained excerpts carry no citation command at all, so no bibliography entry of the article is transcribed and no reference is redistributed. Reference adaptation: none was needed and none was made; every reference inside the delivered unit resolves inside it, so no unresolved reference or citation is produced. Printed slips kept literal: no printed word, symbol, hypothesis or number of the counted statement or of its proof was altered. Layout and class: the article's own class is \texttt{siamart171218} with packages including amsfonts, graphicx, caption, subcaption, mathtools, amssymb, multirow, amsmath, algorithm, algpseudocode, hyperref, array, mathabx, bm; the wrapper is the lane's pinned \texttt{amsart} editorial wrapper at 10pt on A4, which restates the theorem-like environments the retained text needs and adds the article's own macro definitions that the retained text needs (none), with extra wrapper packages (bm, bbm, dsfont, xspace, cancel, mathtools). The delivered numbering is the unit's own. Assets: no retained span contains a figure environment, a graphics-inclusion command, a PDF-page inclusion, a table or a TikZ picture; no image file is copied into this package, the package declares no asset dependency and no image file is redistributed with it. Licence: version arXiv:2504.19308v4 of this article is distributed under the Creative Commons Attribution 4.0 International licence, as recorded in the arXiv licence metadata for this exact version and shown on the abstract page https://arxiv.org/abs/2504.19308v4. A full rights notice is delivered at licenses/arxiv-2504.19308-CC-BY-4.0.txt. Characters: every retained excerpt is pure ASCII, so the delivered engine, XeLaTeX with its default Latin Modern text font, reports no missing character.
\begin{lemma}\label{lem:circulant-conjugation}

Let $R \in \mathbb{C}^{n \times n}$ be circulant, i.e. $R(i,k) = r_{k-i \ (\mathrm{mod}\ n)}.$ Let $D = \mathrm{diag}(\omega^0, \omega^1, \dots, \omega^{n-1})$ with $\omega = e^{-2\pi i / n}.$ Then for any integer $j$, the matrix $\widetilde{R} := D^j R D^{-j}$ is also circulant. Consequently, $D^j R = \widetilde{R} D^j.$
\end{lemma}

\begin{theorem}[Residual bound for Circulant Decomposition]
\label{th:cd-deterministic}
Let $A, B \in \mathbb{R}^{n \times n}$ and let $\Delta A, \Delta B$ be their residuals from a truncated circulant decomposition. Then,
\[
\|\Delta A \Delta B\|_F
\;\le\;
\frac{(\mu_A \mu_B)^{1/4}}{\sqrt{n}}\;
\|\Delta A\|_F \, \|\Delta B\|_F.
\]
\end{theorem}

\begin{proof}
Using the circulant decomposition,
\[
\Delta A \Delta B
=
\sum_{j \notin S^A} \sum_{l \notin S^B}
R_j^A D^j R_l^B D^l.
\]
In the subsequent arguments we use $\sum_{j+l = s}^*$ to denote the indices $j \notin S^A$ and $l \notin S^B$ with their sum as $s$. 

Using the Lemma \ref{lem:circulant-conjugation}, we have $D^j R_l^B = \tilde{R}_l^B D^j$, and we obtain:
\[
\Delta A \Delta B
=
\sum_s R_s D^s,
\quad
R_s := \sum_{j+l=s}^* R_j^A \tilde{R}_l^B.
\]
By orthogonality, $\|\Delta A \Delta B\|_F^2 = \sum_s \|R_s\|_F^2.$ Since circulant matrices are diagonalized by the Fourier transform,
\[
\|R_s\|_F^2 = \sum_{i=1}^n |\lambda_i(R_s)|^2,
\quad
\lambda_i(R_s) = \sum_{j+l=s}^* \lambda_i(R_j^A)\lambda_i(R_l^B).
\]
Let $a_j^{(i)} := \lambda_i(R_j^A),
\quad b_l^{(i)} := \lambda_i(R_l^B).$

Then, $\lambda_i(R_s) = \sum_{j+l=s}^* a_j^{(i)} b_l^{(i)}$, and we have the circular convolution $\lambda_i(R_s) = (a^{(i)} \circledast b^{(i)})(s)$. 

Summing over $s$,
\[
\sum_s |\lambda_i(R_s)|^2
=
\|a^{(i)} \circledast b^{(i)}\|_2^2
\le
\|a^{(i)}\|_2^2 \|b^{(i)}\|_2^2.
\]
Thus,
\[
\sum_s |\lambda_i(R_s)|^2
\le
\left(\sum_{j \notin S^A} |\lambda_i(R_j^A)|^2\right)
\left(\sum_{l \notin S^B} |\lambda_i(R_l^B)|^2\right)
= X_i Y_i.
\]
Summing over, $\|\Delta A \Delta B\|_F^2
\le \sum_i X_i Y_i.$ Applying Cauchy--Schwarz inequality,
\[
\sum_i X_i Y_i
\le
\left(\sum_i X_i^2\right)^{1/2}
\left(\sum_i Y_i^2\right)^{1/2}.
\]

Given $\sum_i X_i = \|\Delta A\|_F^2$, we have
\[
\sum_i X_i^2 \le (\max_i X_i)\sum_i X_i
=
\frac{\mu_A}{n}\|\Delta A\|_F^4,
\]
and similarly for $Y_i$. Therefore,
\[
\|\Delta A \Delta B\|_F^2
\le
\frac{\sqrt{\mu_A \mu_B}}{n}
\|\Delta A\|_F^2 \|\Delta B\|_F^2.
\]

Applying square roots on both sides of the above completes the proof.
\end{proof}

\end{document}
